\clearpage
\SourceLeaf{274}%
\begin{center}
\textbf{[METAPHYSIK DER GEOMETRIE.]}

\medskip
[\textsc{W. Sartorius von Waltershausen.} \emph{Gauss zum Gedächtniss. Leipzig 1856. S. 80---81.}]
\end{center}

\{\mbox{.\,.\,.\,.\,.\,.\,.\,.\,.} In seiner frühsten Jugend habe ihm [\textsc{Gauss}] die Geometrie wenig Interesse eingeflösst, welches sich erst später bei ihm in hohem Masse entwickelt habe.

Es war besonders merkwürdig und überaus lehrreich, von \textsc{Gauss} die Fundamente, auf denen die Mathematik basirt ist, blossgelegt und sie gegen die Metaphysik scharf abgegrenzt zu erblicken. Obgleich er über diese Fragen nie etwas veröffentlicht hat, so steht doch zu vermuthen, dass sich darüber in seinem wissenschaftlichen Nachlass einiges vorfinden wird. In früherer Zeit, als seine Lebensrichtung noch nicht entschieden war und er daran denken musste, dass er vielleicht als Lehrer der Mathematik irgendwo aufzutreten habe, hatte er sich in dieser Aussicht ein Papier ausgearbeitet, welches noch in seinen letzten Jahren vorhanden gewesen sein soll und auf dem er die Anfänge der Mathematik philosophisch entwickelt hatte. Ob dasselbe sich jetzt noch vorfinden wird, ist zweifelhaft.

Die Geometrie betrachtete \textsc{Gauss} nur als ein consequentes Gebäude, nachdem die Parallelentheorie als Axiom an der Spitze zugegeben sei; er sei indess zur Überzeugung gelangt, dass dieser Satz nicht bewiesen werden könne, doch wisse man aus der Erfahrung, z. B. aus den Winkeln des Dreiecks Brocken, Hohehagen, Inselsberg, dass er näherungsweise richtig sei. Wolle man dagegen das genannte Axiom nicht zugeben, so folge daraus eine andere ganz
